\documentclass[11pt]{amsart}
\usepackage{graphicx}
\usepackage{amsfonts,amsmath,amssymb,amsthm}
\usepackage{mathtools}
\usepackage[a4paper,margin=1in]{geometry} 
\usepackage{hyperref}
\usepackage{tikz}
\usepackage{comment}
\usetikzlibrary{positioning,calc}

% TODO
\hbadness=10000

\DeclarePairedDelimiter\norm{\lVert}{\rVert}
\DeclarePairedDelimiter\fracpart{\{}{\}}
\DeclarePairedDelimiter\floor{\lfloor}{\rfloor}
\DeclarePairedDelimiter\ceil{\lceil}{\rceil}

\newcommand{\set}[1]{\left\{#1\right\}}
\newcommand{\abs}[1]{\left|#1\right|}
\newcommand{\paren}[1]{\left(#1\right)}
\newcommand{\Int}[1]{\mathbb Z_{#1}}
\newcommand{\Inttilde}[1]{\tilde{\mathbb Z}_{#1}}
\newcommand{\Unit}[1]{\Int{#1}^\times}

\newcommand{\R}{\mathbb{R}}
\newcommand{\Z}{\mathbb{Z}}
\newcommand{\Zko}{\Int{k+1}}
\newcommand{\Zp}{\Int{p}}
\newcommand{\rbf}{\mathbf{r}}
\newcommand{\abf}{\mathbf{a}}
\newcommand{\vbf}{\mathbf{v}}
\newcommand{\Vbf}{\mathbf{V}}
\newcommand{\ubf}{\mathbf{u}}
\newcommand{\Ubf}{\mathbf{U}}
\newcommand{\wbf}{\mathbf{w}}
\newcommand{\Wbf}{\mathbf{W}}
\newcommand{\nequiv}{\not\equiv}
\newcommand{\onetok}{(1, 2, \dots, k)}

\newtheorem{theorem}{Theorem}[section]
\newtheorem{lemma}[theorem]{Lemma}
\newtheorem{proposition}[theorem]{Proposition}
\newtheorem{corollary}[theorem]{Corollary}
\newtheorem{conjecture}[theorem]{Conjecture}

\theoremstyle{definition}
\newtheorem{definition}[theorem]{Definition}
\newtheorem{example}[theorem]{Example}

\theoremstyle{remark}
\newtheorem{remark}[theorem]{Remark}
\newtheorem{claim}[theorem]{Claim}

\title{Eleven, twelve, thirteen lonely runners}
\author[T. Sungkawichai]{Touch Sungkawichai}
\date{\today}
\address{London, UK}
\email{touchsungkawichai@gmail.com}
\author[T. Trakulthongchai]{Tanupat Trakulthongchai}
\address{St John's College, University of Oxford, Oxford, UK}
\email{tanupat.trakulthongchai@sjc.ox.ac.uk}
\date{\today}
\subjclass[2020]{11K60}

\begin{document}
\begin{abstract}
    Wills stated the Lonely Runner Conjecture: if $k+1$ runners run with distinct and constant speeds around a unit circle from a 
    common starting point, each runner is at least $1/(k+1)$ away from every other runner at some point in time.
    A computational method developed by Rosenfeld and the second author verified the conjecture for up to $10$ runners.
    We further refined this method and provided a computer-assisted proofs of the Lonely Runner Conjecture for up to $13$ runners.
\end{abstract}
\maketitle

\section{Introduction}
The Lonely Runner Conjecture arose in the study of Diophantine approximation \cite{Wills},
though it has been interpreted in the context of view obstruction and billiard trajectories 
\cite{cusick1984view}, covering radii of zonotopes \cite{henze_covering_2017}, and the chromatic number of a 
distance graph \cite{zhu2002circular}. 
It was conjectured independently by Wills \cite{Wills} in 1967 and Cusick \cite{cusick1984view} in 1984. 
In the most popular form, the conjecture asserts that if $k+1$ runners run with distinct,
constant speeds around a unit circle, each of them will be \emph{lonely} at some time. 
Here, a runner is said to be lonely if every other runner is at a distance at least $\tfrac{1}{k+1}$ from them along the circle.

\begin{conjecture} \label{lonelyconj}
  Let $u_1,\ldots,u_{k+1}$ be distinct real numbers. For each $i\in\{1,\ldots,k+1\}$, there is a real number $t$ such that
  \[
    \norm{t u_i-tu_j} \ge \frac{1}{k+1} \quad \text{for all } i \neq j.
  \]
\end{conjecture}

The conjecture at each specific value of $k$ is called the Lonely Runner Conjecture for $k+1$ runners.

The Lonely Runner Conjecture is currently known to hold for cases of up to 10 runners.
The two smallest cases for $k$ are trivial. 
Betke~and~Wills~\cite{fourrunners} gave the proof for $k=3$ (the 4 runners case), 
and Cusick~and~Pomerance~\cite{cusick1984view} verified the conjecture for $k=4$ (the 5 runners case).
In the same paper that proved the case $k=5$, Bohman,~Holzman,~and~Kleitman~\cite{Bohman} showed that it suffices 
to consider integer speeds, transforming the problem from a continuous one to a discrete one. 
This reduction, combined with fixing the speed of one runner to $0$, produces the following
equivalent statement to Conjecture~\ref{lonelyconj}.

\begin{conjecture}
    Any $\ubf\in\Z_{>0}^k$ has the \emph{LR property}, that is, there is a real number $t$ such that
    \[\norm{tu_i}\geq\frac{1}{k+1} \quad \text{for all } i = 1,\dots,k.\] 
    
    We denote this statement by $LRC(k)$.
\end{conjecture}

To match the conjecture's settings, we refer to the vector $\ubf = (u_1,\ldots,u_k)$ as a \emph{speed set}.

Before 2025, the conjecture was known to be true for $k\leq6$, where the final case was proven by 
Barajas~and~Serra~\cite{barajas_lonely_2008} in 2008. In September 2025, Rosenfeld~\cite{Rosenfeld} introduced a computational 
method that settled $k=7$ (and later $k=8$ \cite{rosenfeld_lonely_2026}). In doing so, he relies on the finite-checking
result of Malikiosis, Santos, and Schymura~\cite{Malikiosis}, which follows from similar results of Tao~\cite{Tao} and
Giri~and~Kravitz~\cite{giri_structure_2026}. 
The~second~author~\cite{Trakulthongchai} subsequently refined this method to prove the conjecture for the case of $k\in\{8,9\}$ in November 2025.

In this paper, we further develop this computational method to verify the Lonely Runner Conjecture for $k\in\{10,11,12\}$. 
Our overall strategy is outlined in the next section within the context of the computational method of \cite{Rosenfeld} 
and \cite{Trakulthongchai}. 
Sections~\ref{sec:quotient}, \ref{sec:sieves} and \ref{sec:onetok} explain the techniques that help simplify verification.
Section~\ref{sec:algo} describes the algorithms used for each $k\in\{10,11,12\}$, 
the results and implementation of which are presented in Sections~\ref{sec:result} and \ref{sec:implement}, respectively. 
% Lastly, we discuss in Section~\ref{sec:further} how this method may be further extended. 

The article adopts the following choices for notation: $\Int{n}$ denotes the additive group modulo $n$,
$\Unit{n}$ is the multiplicative group modulo $n$, and $\Int{n,l} := \Int{ln} \setminus n\Int{l}$ is a subset of $\Int{ln}$.
Moreover, let $\pi_{m \to n}:\Int{m} \to \Int{n}$ denote the natural projection map when $n\mid m$
and $\pi_n: \Z \to \Int{n}$ denote the projection modulo $n$.
The application of a scalar function to vectors is coordinate-wise. $\widehat{\cdot}$ indicates the omission of a 
certain coordinate in a vector. The fractional part of $x\in\R$ is denoted by $\fracpart{x}$.

\section{Overview}
We follow the computational framework laid by Rosenfeld~\cite{Rosenfeld} and adopt the terminology\footnote{
  Compared to \cite{Trakulthongchai}, the main difference is replacing $I(k,l,p)$ with $I(k,p,l)$.
} and techniques from \cite{Trakulthongchai}.

We begin by recalling the key result for the current finite-verification framework by Rosenfeld, 
building from a result obtained by Malikiosis,~Santos,~and~Schymura~\cite{Malikiosis}:
\begin{lemma}[Corollary 3, \cite{Rosenfeld}] \label{lem:proddivB}
    If $LRC(k-1)$ holds, then any counterexample $\ubf$ to $LRC(k)$ must have
    \[ u_1\cdots u_k < B_k \quad \text{where} \quad B_k := \paren{\frac{\binom{k+1}{2}^{k-1}}{k}}^k. \]
\end{lemma}

This implies that $LRC(k)$ can be verified by checking all positive integral speed sets such that $u_1 \cdots u_k < B_k$.
Rosenfeld approached this bound by accumulating a collection of small proofs,
each asserting that for a particular prime $p$, any counterexample $\ubf$ must satisfy $p \mid u_1 \dots u_k$.
By collecting enough small proofs for a set of prime $P$ such that $\prod_{p \in P} > B_k$, it follows that 
no counterexample can exist\footnote{
  Each small proof not need be prime if the product was replaced by lcm. However, we always use prime moduli in this article.
}.
These small proofs are obtained computationally with the help of the following Lemma~\ref{lem:divp}. 
We first restate the definition of \emph{improper} sets proposed in \cite{Trakulthongchai}.

\begin{definition}[Definition 3.1, \cite{Trakulthongchai}] \label{def:improper}
    Let $p$ be a prime and $k,l$ be positive integers.
    A speed set $\vbf\in\Int{p,l}^k$ is $(k,p,l)$-\emph{improper} if both conditions hold: 
  \begin{itemize}
    \item for any $t \in \Int{lp}$ there exists $i$ such that $\norm{\frac{tv_i}{lp}} < \frac{1}{k+1}$.
    \item for all $i$, $\gcd(lp, v_1, \dots, \widehat{v_{i}}, \dots, v_k) = 1$.
  \end{itemize}
   If $\vbf$ is not improper, it is \emph{proper}. Moreover, let $I(k,p,l)$ denote the set of all improper speed sets.
\end{definition}

By the definition, it is clear that permuting coordinates preserves both conditions. 
In other words, $I(k, p, l)$ carries a natural $S_k$ symmetry on $I(k, p, l)$ which both 
\cite{Rosenfeld} and \cite{Trakulthongchai} exploited in their search. 
Moreover, it can be seen that replacing $v_i$ with $-v_i$ at any coordinate also preserves both conditions. 
We suppress this detail but note that equality of tuples understood throughout is modulo these symmetries.

One of Rosenfeld's \cite{Rosenfeld} key insights is to restrict the time parameter $t$ to be a rational number with denominator $lp$.
By doing so, the time point $t$ witnessing loneliness of a speed set $\vbf \in \Int{lp}^k$ applies to all speed set $\ubf$ in the 
preimage $\pi^{-1}_{lp}(\vbf)$.

Our computer-assisted proofs follow the method of~\cite{Trakulthongchai}, which establishes the required small proofs 
by verifying emptyness of the set $I$ at certain parameters; a direct result from the following lemma.

\begin{lemma}[Lemma 3.3, \cite{Trakulthongchai}] \label{lem:divp}
    For $k\geq3$, if $LRC(k-1)$ holds and $I(k,p,l)$ is empty for some $l$, then any counterexample
    $\ubf$ to $LRC(k)$ satisfies $p \mid u_1 \cdots u_k$. 
\end{lemma}

The main advance of \cite{Trakulthongchai} is the observation that allows $I(k,p,cl)$ to be computed efficiently.
The following lemma enables the \emph{lifting} process of $I(k,p,l)$ to $I(k,p,cl)$ which is more fine-grained.
This lifting is done in hope that the extra structure of $\Int{clp}$ will help filtering out speed sets that cannot 
be counterexamples to the conjecture.

\begin{lemma}[Lemma 3.5, \cite{Trakulthongchai}] \label{lem:shadowing}
  For any $k, p, l, c$, the inclusion 
  \[ (\pi_{clp \to lp}) I(k,p,cl) \subseteq I(k,p,l) \]
  holds.
\end{lemma}

As seen the previous lemmas, the central goal is to show that $I(k, p, l) = \emptyset$ for some value $l$. 
This motivates the following definition of an \emph{eventually proper} speed set, 
which is the main subject of interest in this paper.
\begin{definition}
  \label{def:eventualproper}
  A speed set $\vbf \in \Int{p,1}^k$ is said to be \emph{eventually $(k,p)$-proper} if there is $l$ such 
  that \[ \vbf \notin (\pi_{lp \to p}) I(k, p, l). \]
\end{definition}

The concept of eventual properness simplifies the discussion and implementation; while many speed sets in $\Int{p,1}^k$ are not proper,
it is enough that they are all eventually proper. 

\begin{lemma} \label{lem:evenpropdivp}
  Let $k \ge 3$ and assume that $LRC(k-1)$ is true. If every $\vbf \in \Int{p,1}^k$ is eventually $(k,p)$-proper,
  then any counterexample $\ubf$ to $LRC(k)$ has $p \mid u_1\cdots u_k$.
\end{lemma}
\begin{proof}
  Since every $\vbf\in\Int{p,1}^k$ is eventually proper, for each $\vbf$ there is an integer $l_\vbf$ for which 
  \[ \vbf \notin (\pi_{l_\vbf p \to p}) I(k, p, l_\vbf). \]
  From there, pick $l$ to be the least common multiple of all $l_\vbf$ where $\vbf\in\Int{p,1}^k$.
  By Lemma~\ref{lem:shadowing}, 
  \[ (\pi_{l_\vbf p \to p}) (\pi_{lp \to l_\vbf p}) I(k, p, lp) \subseteq (\pi_{l_\vbf p \to p}) I(k, p, l_\vbf p). \]
  Hence as $\vbf$ is not in the set on the right hand side, it follows that, $\vbf \notin \pi_{lp \to p} I(k, p, lp)$.
  As this applies to all $\vbf$, it follows that $I(k, p, lp)$ is empty. Lemma~\ref{lem:divp} yields the desired conclusion. 
\end{proof}

\section{
  Quotient modulo \texorpdfstring{$p$}{p}
} \label{sec:quotient}
The framework introduced in \cite{Trakulthongchai} has a bottleneck in the first phase of computation, namely,
the calculation of $I(k, p, 1)$. In this section, we aim to exploit the $\Unit{p}$ symmetry of $I(k,p,l)$ to reduce 
the amount of computation needed for $I(k, p, 1)$ and its subsequent lifts by a factor of $p-1$, which 
in turn reduce the amount of computation work needed for $I(k, p, 1)$. 
In our implementation, accounting for the interplay with the existing symmetries, 
this yields an overall reduction in computation by a factor of $\Theta(p/k)$.

For a fixed prime $p$, consider the relation $\sim$ in $\Int{p,l}^k$ defined as
\[ \ubf \sim \vbf \iff \ubf = a\vbf \quad \text{for some } a\in\Unit{p}.\] 
It follows that $\sim$ is an equivalence relation. More importantly, this relation also preserves properness.

\begin{lemma} \label{lem:simpreserve}
Let $p$ be prime and $l\ge1$ and let $\ubf,\vbf$ be speed sets.
\begin{enumerate}
  \item If $\ubf \sim \vbf$, then $\ubf$ is improper if and only if $\vbf$ is improper.
  \item For $\ubf,\vbf\in\Int{p,1}^k$, if $\ubf \sim \vbf$, then $\ubf$ is eventually proper if and only if $\vbf$ is eventually proper.
\end{enumerate}
\end{lemma}

\begin{proof}
  Let $\ubf = a\vbf$ with $a \in \Unit{p}$.
  \begin{enumerate}
    \item As $a$ is a unit, $a^{-1} \ubf = \vbf$ so, 
    by symmetry, it suffices to show that improperness of $\vbf$ implies improperness of $\ubf$.
    Since $\gcd(a,p)=1$, it follows that  
    \[\gcd(lp, u_1, \dots, \widehat{u_i}, \dots, u_k) = \gcd(lp,av_1,\ldots,\widehat{av_{i}},\ldots,v_k) = \gcd(lp, v_1, \dots, \widehat{v_i}, \dots, v_k) = 1,\]
    thus the gcd condition in Definition~\ref{def:improper} is preserved.
    Moreover, for any $t \in \Int{lp}$,
    \[ \norm*{\frac{t u_i}{p}} = \norm*{\frac{t(av_i)}{p}} = \norm*{\frac{(at)v_i}{p}}, \]
    and since $at \in \Int{lp}$, the distance condition holds for $\ubf$ if and only if it holds for $\vbf$.
    Thus $\ubf$ is improper if and only if $\vbf$ is improper.

    \item Again, by symmetry, it suffices to show that $\ubf$ is eventually proper assuming $\vbf$ is.
    For brevity, let $\pi = \pi_{lp\to p}$.
    As $\vbf$ is eventually proper, there exists $l$ such that $\vbf \not\in \pi I(k, p, l)$.
    In other words, $\pi^{-1}(\vbf) \cap I(k, p, l) = \emptyset$.
    Since the restriction $\pi^{-1} : \Unit{lp} \to \Unit{p}$ is surjective by the Chinese Remainder Theorem, 
    we may choose $b \in \Unit{lp}$ with $\pi(b) = a^{-1}$.

    For any $\ubf' \in \pi^{-1}(\ubf)$ we let $\vbf' = b\ubf'$. Then, 
    \[
      \pi(\vbf') = \pi(b)\pi(\ubf') = a^{-1}\ubf = \vbf,
    \]
    so $\vbf' \in \pi^{-1}(\vbf)$ and hence $\vbf' \notin I(k, p, l)$, so it is proper.
    By part (1), since $\vbf' \sim \ubf'$, $\ubf'$ is also proper.
    Thus $\pi^{-1}(\ubf)$ contains no improper speed set, ensuring that $\ubf$ is eventually proper.
  \end{enumerate}

  Therefore, both statements hold.
\end{proof}

Lemma~\ref{lem:simpreserve} essentially says that eventual properness only need to be checked for a single representative 
from each equivalence class of $\sim$. Therefore, we write $\tilde{I}(k,p,l)$ and $\Inttilde{p,l}^k$
to refer to the quotient of $I(k, p, l)$ and $\Int{p,1}^k$ by $\sim$ respectively. 
By a standard abuse of notation, we identify each quotient with a fixed choice of representative speed 
sets taken from $I(k, p, l)$ or $\Int{p,1}^k$. With these new notations, we state a key result of this section.

\begin{corollary} \label{cor:equivclasscheck}
  Let $k \ge 3$ and assume that $LRC(k-1)$ is true. If every $\vbf \in \Inttilde{p,1}^k$ is eventually $(k,p)$-proper,
  then any counterexample $\ubf$ to $LRC(k)$ has $p \mid u_1\cdots u_k$.
\end{corollary}
\begin{proof}
    The statement follows immediately from Lemma~\ref{lem:evenpropdivp} and part (2) of Lemma~\ref{lem:simpreserve}.
\end{proof}

\section{Intermediate sieves} \label{sec:sieves}

Recall that our ultimate goal is to show the eventual emptiness of $\tilde{I}$.
In this section we describe two techniques that substantially reduce the cost of this procedure in practice.

As we can identify $\tilde{I}(k, p, l)$ as a subset of $I(k, p, l)$, many results from \cite{Trakulthongchai} are readily available. 
For example, lifting $\tilde{I}(k, p, l)$ to $\tilde{I}(k, p, cl)$ can be thought in the exact same way as lifting 
$I(k, p, l)$ to $I(k, p, cl)$.

\subsection{Efficient Lifting}
The cost of lifting $\tilde{I}(k,p,l)$ to $\tilde{I}(k,p,cl)$ is proportional to \[ |\tilde{I}(k,p,l)| \cdot c^k\] 
since each improper set modulo $lp$ has $c^k$ candidates (preimage modulo $clp$) that need verification.
When $k$ is large, the value of $c$ has a greater impact on computation time. For example, for $k=11$, 
the $c=3$ lifts are almost 90 times more expensive than those with $c=2$.

It is essential to carefully choose the order of each lift.
Since the $c=2$ lifts are cheap, performing them multiple times can help if they reduce the size of the search space,
i.e., if $|\tilde{I}(k, p, 2l)| < |\tilde{I}(k,p,l)|$.
Thus, employing a number of cheap unnecessary lifts before performing the expensive mandatory lifts, for example, 
adding $c=2$ lifts before $c=11$ lifts when $k=10$, can shorten the overall computation time 
by reducing the amount of computation needed in larger moduli.

\subsection{Backward projection}

Just as lifting can reduce the size of the search space, projection can achieve a similar effect.
After computing $\tilde{I}(k, p, l)$, projecting the speed sets back into $\Int{p}$ via the map $\pi_{pl \to p}$
can collapse several improper speed sets from $\tilde{I}(k, p, l)$ to a representative by maintaining only the structure 
in $\Int{p}$, discarding all other structures of $\Int{lp}$.
This action loses fine-grained detail of each improper speed set, but reduces the number of improper speed sets that need calculation.
Notice that $\abs{\pi_{pl \to p} S} \le \abs{S}$.

Crucially, this operation is lossless with respect to eventual properness
because if a speed set $\vbf$ is not in the projected $\pi_{lp \to p} \tilde{I}(k, p, l)$, then by 
Definition~\ref{def:eventualproper} and Lemma~\ref{lem:simpreserve}, it follows that $\vbf$ is eventually proper.
This fact implies that taking projection of a search space at any point will not skip over any speed sets.

\section{Eventual properness of
  \texorpdfstring{$(1, 2,\ldots, k)$}{(1, 2, ..., k)}
} \label{sec:onetok}

In the case of $k\in\{10,12\}$, where $k+1$ is prime, the cost of lifting with $c=k+1$ is enormous, so we wish to bypass this lift.
Employing the techniques mentioned in Section~\ref{sec:sieves} led to a discovery that 
after lifting $\tilde{I}(10,p,1)$ with $c=2,3$ for a few times and projecting back, 
the only representative speed set in $\tilde{I}(10,p,1)$ that cannot be shown to be eventually proper is the speed set 
$( 1, 2, \ldots, 10 )$. This holds for many primes and the analogous observation holds for many primes in the case of 
$k=12$ as well.

It is unsurprising that $\onetok$ remains as it can never obey the gcd condition and it is only lonely at time 
$t = \tfrac{s}{k+1}$ for $s \in \Z$. Consequently, there is no time of the 
form $t = \frac{r}{pl}$ where $r \in \Z$ that would satisfy the condition; a lift with $k+1 \nmid c$ will not resolve these.
What is perhaps more surprising is that no other speed sets survive the pruning.

The arguments in this section aim to exploit the structure of the problem when $k+1$ is prime. 
With additional care in tracking divisibility conditions, most of the results can be extended
to the case where $k+1$ is not prime, but the results are necessarily weaker.

The key structural observation is that the speed set $\onetok$ modulo $p$ projects naturally onto $\Zko$. 
Instead of considering the preimage of $\onetok$ directly, we isolate the effect of $p$
by considering time points of the form $t=\frac{q}{(k+1)p}$ as a partial fraction $t=\frac{s}{k+1}+\frac{r}{p}$. 
We first define a \emph{resolvability} of speed set modulo $\Zko$ tailored to the speed set $\onetok$.

\begin{definition}
\label{def:resolvable}
  Let $R_n(x)=\floor{ n \fracpart{x} }$.
  Equivalently, $R_n(x)$ is the integer $k$ such that $\fracpart{x}$ lies in the $1/n$-arc $[k/n, (k+1)/n)$.

  Let $\vbf \in \Zko^k$. 
  We write \[ \rbf(x)=\left(R_{k+1}(x), \; R_{k+1}(2x), \; \ldots, \; R_{k+1}(kx)\right) \] for $x\in\R$ and 
  say that $\vbf$ is $(k,p)$-\emph{resolvable} if any of the following holds:
  \begin{itemize}
      \item (scattered) $\vbf$ contains no zero coordinate
      \item (divisible) $\vbf$ contains at least $k/2$ zero coordinates
      \item (lonely) there is $s \in \Zko$ and $r \in \Z$ for which
      \[ s\vbf + \rbf(r/p) \in \{1,\ldots,k-1\}^k\] when considered modulo $k+1$.
  \end{itemize}
\end{definition}

Notice that $\rbf$ has a periodicity of $1$, thus it suffices to consider only those $r \le p$. 
Moreover, it can be seen that all dependence on $p$ is eliminated except for that arising from $\rbf(\frac{1}{p} \Zp)$
in the lonely condition. 

In this section, we show that, in a sense, it is enough to consider the projected $\onetok$ in $\Zko$
once and the result can be reused in larger moduli.

We first recall a variation of a folklore result known by many as evidenced in the Lonely Runner literature.
\begin{lemma}
\label{folklore}
    Suppose that $LRC(n)$ holds for all $n<k$. Let $p$ be a prime that divides $m$ speeds in $\vbf$, 
    where $\ceil*{ \frac{2p}{k+1} } <\frac{p}{k-m}$, and $\gcd(v_1,\ldots,v_k)=1$.
    Then $\vbf$ has the LR property.
\end{lemma}
\begin{proof}
  We can assume that
  $v_1,\ldots,v_m$ are divisible by $p$ and $v_{m+1},\ldots,v_k$ are not. 
  Since $m<k$, $LRC(m)$ holds and there is a time $T \in \R$ when $\norm{Tv_i} \ge \frac{1}{m+1}$ for all $i\leq m$.
  
  By considering time of the form $t=T+\frac{r}{p}$ for $r \in \Zp$, we have that 
  \[ \norm{tv_i} = \norm{Tv_i} \ge \frac{1}{m+1}>\frac{1}{k+1} \quad \text{for all } i \le m \] as $p$ divides $v_i$. 
  For $i > m$, on the other hand, 
  \[ \set{ \fracpart*{\paren{T + \frac{r}{p}} v_i} : r \in \Zp } = \set{ \fracpart*{Tv_i + \frac{r}{p}} : r \in \Zp } \]
  as multiplication by $v_i$ permutes elements in $\Zp$.

  Clearly, the latter set is the set of $p$ evenly spaced points on the unit circle, 
  so at most $\ceil*{\frac{2}{k+1}\cdot p }$ points can intersect with the interval $\paren{-\frac{1}{k+1},\frac{1}{k+1}}$. 
  In other words, at most $\ceil*{ \frac{2p}{k+1} }$ points are not lonely for each $i\in\{m+1,\ldots,k\}$. 
    
  Since there are $k-m$ such indices, at most $(k-m) \ceil*{ \frac{2p}{k+1} }$ choices of $r\in\Zp$ are eliminated. 
  From $(k-m) \ceil*{ \frac{2p}{k+1} } < p$, there must exist $r\in\Int{p}$ that makes $\norm{tv_i}\geq\frac{1}{k+1}$
  for all $i\geq m+1$.
\end{proof}

With the lemma, we can show that resolvability of all speed sets entails eventual properness of $\onetok$.
\begin{lemma}
\label{lem:resolvable}
    Suppose that $LRC(n)$ holds for all $n<k$. Assume that all speed sets are $(k, p)$-resolvable and $k+1$ is prime.
    Then $\onetok$ is eventually $(k,p)$-proper.    
\end{lemma}
\begin{proof}
    Let us denote $\vbf = \onetok$ as a speed set in $\Int{p, 1}^k$.
    Let $\ubf \in \Int{p(k+1)}^k$ be a preimage of $\vbf$ by $\pi_{p(k+1) \to p}$, 
    that is for each $i$, $u_i = i + n_ip$ for some $0 \le n_i \le k$.
    By considering time of the form $t=\frac{s}{k+1}+\frac{r}{p}$ for $s\in\Z$ and $r\in\Z$, 
    the fractional part of $tu_i$ is given by
    \[ \fracpart{tu_i}=\fracpart*{\left(\frac{s}{k+1}+\frac{r}{p}\right) u_i}=\fracpart*{\frac{su'_i}{k+1}+\frac{ri}{p}} \] 
    where $\ubf' := (u'_1, \dots, u'_k) = \pi_{k+1} \ubf$.

    By assumption, $\ubf'$ is resolvable, thus one of the following cases holds.
    \begin{itemize}
        \item $\ubf'$ is scattered. In this case, pick $r=0$ and $s=1$ and we have 
          \[ \fracpart{tu_i}=\fracpart*{\frac{u'_i}{k+1}} \]
          Since $u_{\tilde{I}} \ne 0$, $\norm{tu_i} \ge \frac{1}{k+1}$ for all $i$.
          Hence $\ubf$ has the LR property.
        \item $\ubf'$ is divisible. Assume that $\ubf'$ contains $m \ge k/2$ zero coordinates. 
          This precisely means that there are exactly $m$ coordinates of $\ubf$ that are divisible by $k+1$.
          By Lemma~\ref{folklore} with $p=k+1$ and $m \ge k/2$ 
          (note that the inequality
          \[\ceil*{ \frac{2p}{k+1} } = 2 < \frac{2(k+1)}{k-1} \leq \frac{p}{k-m}\] 
          holds),
          $\ubf$ has the LR property.
        \item $\ubf'$ is lonely.
          If this is the case, there is a choice of $r$ and $s$ such that 
          \[ s\ubf' + \rbf(r/p) \in \{1,\ldots,k-1\}^k \]
          when considered modulo $k+1$. 
          This means that for each coordinate $i$, there are $A_i$ and $B_i$ such that 
          $\fracpart{\frac{sv_i}{k+1}}=\frac{A_i}{k+1}$ and 
          $\fracpart{\frac{ri}{p}}\in\left[\frac{B_i}{k+1},\frac{B_i+1}{k+1}\right)$ where
          where $C_i:=A_i+B_i$ is in $\set{1,\ldots,k-1}$ when considered modulo $k+1$. 
          Therefore,
          \[\fracpart{tu_i}=\fracpart*{\frac{su_i}{k+1}+\frac{ri}{p}}\in\left[\frac{C_i}{k+1},\frac{C_i+1}{k+1}\right)\]
          and in particular the set on the right hand side does not intersect with $[0,\frac{1}{k+1}) \cup (\frac{k}{k+1},1)$.  
          Since this holds for every $i$, $\ubf$ has the LR property.
    \end{itemize}
    
    As all $\ubf$ in the preimage of $\vbf$ have the LR property, $\vbf$ must be eventually proper.
\end{proof}

Still, the most demanding task in Lemma~\ref{lem:resolvable} is verifying that all speed sets are
$(k, p)$-resolvable for all $p$\footnote{
  This can be done for $k=10$; it took 3 hours to check all necessary $p$.
  It is much less feasible for $k=12$ where we estimate a few months of constant running.
}, because it involves checking $\Theta((k+1)^k)$ speed sets. 
Fortunately, this verification is unnecessary when $p$ is large. In order to show that, we first need the following proprosition.

\begin{proposition}
\label{prop:nullstellensatz}
    Let $p=k+1$ be prime and $p\geq3$. If the vector $\vbf\in\Int{p}^k$ has the property that, for every $s, r\in\Unit{p}$,
    \[s\vbf + r\onetok  \notin  \{1,\ldots,p-2\}^k,\]
    and $\vbf$ has at least one zero coordinate, then  $\vbf = \mathbf{0}$.
\end{proposition}
\begin{proof}
    In this proof, we work in the field $\Int{p}$.
    Define \[V_m = \set{v_i+mi:i\in \set{1,\ldots,k} }\] for $m \in \Unit{p}$.
    We say $m\in \Unit{p}$ is \emph{good} if $0 \notin V_m$ and \emph{bad} otherwise. Let $G$ be the set of good $m$'s.

    Now, as $\vbf$ contains at least one zero coordinate, it contains at most $k-1$ non-zero coordinates.
    For each non-zero coordinate $v_i$, there is exactly one choice of $m\in\Unit{p}$ such that $v_i + m i = 0$, whereas for each
    zero coordinate, there cannot be such $m$. Thus, there are at most $k-1$ bad units.
    This implies that $\abs{G} \ge 1$.

    Let $m$ be good. Since $\vbf$ has the property, for any $s\in\Unit{p}$,
    we may choose $r := sm$.
    It follows that 
    \[ s\vbf + r\onetok \not\in \set{1, \dots, p-2}^k \]
    thus there is an index $i$ such that $s (v_i + m i) \in \set{0, -1}$.
    As $m$ is good and $s$ is a unit, it is impossible for the product to be zero, and so it must be $-1$. Hence, 
    \[ v_i + m i = -s^{-1}. \]
    As $s$ was chosen arbitrarily, $V_m = \Unit{p}$.

    From $V_m = \Unit{p}$ when $m$ is good, we have that
    \[k! = \prod_{i=1}^k(v_i+mi) = \prod_{i=1}^k i \cdot \prod_{i=1}^k (v_i i^{-1} + m) = k!\prod_{i=1}^k (v_i i^{-1}+m),\]
    which gives $\prod_{i=1}^k (v_i i^{-1} + m) = 1$.
    On the other hand, when $m$ is bad, then there is an index $i$ such that $v_i + m i = 0$, hence 
    \[ \prod_{i=1}^k (v_i i^{-1} + m) = \prod_{i=1}^k (v_i + m i) i = 0. \]
    Therefore, the polynomial \[ P(X) := \prod_{i=1}^k (v_ii^{-1}+X) \] is an indicator function of $G$ on $\Unit{p}$.
    Note also that since $\vbf$ contains at least one zero coordinate, $P(0) = 0$.

    As $p$ is prime, Fermat's Little Theorem gives another indicator function of $G$, which is
    \[ Q(X) := \sum_{m \in G} (1-(X-m)^k), \] 
    where the sum is not empty as $\abs{G} \ge 1$.
    It also follows that $Q$ is a polynomial of degree $k$.
    
    As $P(0) = 0 = Q(0)$, $P$ and $Q$ are polynomials of degree $k$ that agree on $k+1$ points, hence $P = Q$.
    In particular, their leading coefficients must be equal, that is 
    $1 \equiv -\abs{G} \pmod{p}$. Since $\abs{G} \le k$ as an integer, we must have that $\abs{G} = k$.
    Therefore, there is no bad $m$. Specifically, for any $i\in\{1,\ldots,k\}$, $v_i+mi \neq 0$
    for all $m\in\Unit{p}$. 
    However, this implies $v_i = 0$ for all $i$, thus $\vbf = \mathbf{0}$.
\end{proof}

Proposition~\ref{prop:nullstellensatz} essentially shows that all speedsets are $(k, k+1)$-resolvable. 
The detail can be seen in the following lemma.

\begin{lemma}
\label{lem:resolvekplusone}
    Suppose that $LRC(n)$ holds for all $n<k$ and let $k+1\geq3$ be prime, then $\onetok$ is eventually $(k, k+1)$-proper.
\end{lemma}
\begin{proof}
  By Lemma~\ref{lem:resolvable}, the conclusion follows if every $\vbf \in \Zko^k$ is $(k, k+1)$-resolvable.

  Let $\vbf$ be any speed set in $\Zko^k$. If it is scattered or divisible, then it is resolvable. 
  Otherwise, $\vbf$ contains at least 1 zero coordinate but $\vbf \ne \mathbf{0}$.
  By the contrapositive of Proposition~\ref{prop:nullstellensatz},
  there exist $s,r\in\Unit{k+1}\subseteq\Int{k+1}$ such that
  \[s\vbf + r\onetok  \in \{1,\ldots,k-1\}^k.\] Recall that \[\begin{split}
      \rbf\left(\tfrac{r}{k+1}\right) 
      &=\left(\floor*{ (k+1)\fracpart*{\tfrac{r}{k+1}} },\floor*{ (k+1)\fracpart*{\tfrac{2r}{k+1}} },\ldots, 
        \floor*{ (k+1)\fracpart*{\tfrac{kr}{k+1}}} \right)\\
      &=(r,2r,\ldots,kr)\\&=r(1,2,\ldots,k)
  \end{split}\]
  modulo $k+1$. Therefore,
  \[s\vbf+\rbf \paren{\tfrac{r}{k+1}} \in \{1,\ldots,k-1\}^k,\]thus $\vbf$ is $(k, k+1)$-lonely, which means 
  that it is $(k, k+1)$-resolvable.
\end{proof}

Lemma~\ref{lem:resolvekplusone} establishes eventual properness of $\onetok$ for a single prime, $k+1$. 
However, as noted that all dependence of $p$ lies only in the structure of $\rbf(\tfrac{1}{p} \Z)$, we can formulate 
the following lemma. 

\begin{lemma}
  \label{lem:subsetresolve}
  If $\rbf (\frac{1}{q} \Z) \subseteq \rbf(\frac{1}{p} \Z)$ and $\vbf$ is $(k, q)$-resolvable, 
  then $\vbf$ is also $(k, p)$-resolvable.
\end{lemma}
\begin{proof}
   Note that the only condition that depends on $p$ in Definition~\ref{def:resolvable} is the lonely condition. 
   Suppose that $\vbf$ satisfies the lonely condition for $(k, q)$,
   namely there is a choice of $s \in \Zko$ and $r \in \Z$ such that $s\vbf + \rbf(r/q) \in \{1,\ldots,k-1\}^k$. 
   From $\rbf(\frac{1}{q} \Z_q) \subseteq \rbf(\frac{1}{p} \Zp)$,
   there is $r' \in \Z$ such that $\rbf(\frac{r'}{p}) = \rbf(\frac{r}{q})$ due to the assumed inclusion and hence
   \[s\vbf + \rbf(r'/p) \in \{1,\ldots,k-1\}^k.\] 
   Therefore, $\vbf$ satisfies the lonely condition in the case $(k,p)$.
\end{proof}

Since $\rbf$ has periodicity of $1$, in practice, it enough to check that $\rbf(\frac{1}{q} \Z_{q}) \subset \rbf(\frac{1}{p} \Z_p)$, 
which can be done computationally given primes $p$ and $q$.
The following key lemma combines this fact with the previously proven facts about $\onetok$ to 
remove the need for computational check for any large enough primes.

\begin{lemma}
    \label{lem:corlarge}
    Suppose that $LRC(n)$ holds for all $n<k$ and let $k+1\geq3$ be prime.
    Then $\onetok$ is eventually $(k,p)$-proper for all primes $p\geq k(k+1)$.
\end{lemma}
\begin{proof}
  By Lemma~\ref{lem:subsetresolve}, \ref{lem:resolvekplusone}, and \ref{lem:resolvable}, 
  it suffices to show that $\rbf(\tfrac{1}{k+1} \Z) \subset \rbf(\tfrac{1}{p} \Z)$ for all $p \ge k(k+1)$.
  Fix $p$ to be a prime greater than $k(k+1)$.

  Notice that $\rbf$ is a piecewise constant function with many discontinuity points. In fact, 
  $x$ is a discontinuity of $\rbf$ if and only if it is a discontinuity of $R_{k+1}(ix)$ for some $i$. Those discontinuity 
  points are precisely points in the form of $\frac{a}{i(k+1)}$ for integer $a$ and $1 \le i \le k$.

  Choose any integer $n$ and denote $x := \tfrac{n}{k+1}$. Observe that $x$ is a discontinuity point of $\rbf$ and let $y$ be 
  the smallest number greater than $x$ that is also a discontinuity point of $\rbf$. We know that $y$ can be written as 
  $\tfrac{a}{i(k+1)}$ for some integer $a$ and $1 \le i \le k$, so it follows that 
  \[ y - x = \frac{a}{i(k+1)} - \frac{n}{k+1} = \frac{a - ni}{i(k+1)} \ge \frac{1}{i(k+1)} \ge \frac{1}{k(k+1)} \]

  As $y$ was chosen to be the smallest such discontinuity, $\rbf$ is constant in the interval $[x, y)$. Since the width of said interval
  is at least $\tfrac{1}{k(k+1)} \ge \tfrac{1}{p}$, there must be an integer $m$ such that $\tfrac{m}{p} \in [x, y)$.
  Hence, for such value $m$ 
  \[ \rbf \paren{\frac{n}{k+1}} = \rbf \paren{\frac{m}{p}} \]

  As $n$ was chosen arbitrarily, the desired inclusion follows.
\end{proof}


% -- I don't think this adds anything useful
% The quantity $k(k+1)\approx k^2$ is very reasonable and useful in practice. 
% Without Lemma~\ref{lem:resolvekplusone}, via a similar consideration of the length of gaps 
% between discontinuity points
% the same conclusion to Lemma~\ref{lem:corlarge} exists for $p \ge k^3 - k$ provided that there is some prime $q$ such that all speed sets are $(k, q)$-resolvable. 
% However, this is less useful empirically. 

\section{Algorithms for verification}
\label{sec:algo}
From Corollary~\ref{cor:equivclasscheck}, to show that $p\mid u_1\cdots u_k$ for any counterexample $\ubf$ to $LRC(k)$, 
it suffices to verify that every speed set in $\Inttilde{p,1}^k$ is eventually proper.

For each $k\in\{10,11,12\}$, the verification algorithm consists of two phases: computing $\tilde{I}(k, p, 1)$ and lifting.
The first phase follows the method used in \cite{Trakulthongchai} except deduplicating the structure of $\sim$ as described in 
Section~\ref{sec:quotient}.
We use a lifting diagram to explain the second phase. In a lifting diagram, the first set $S_1$ is always $\tilde{I}(k,p,1)$. 
Each arrow $X_l\rightarrow X_{cl}$ indicates lifting with a constant $c$ and arrow $X_l\xrightarrow{\pi} Y_1$ indicates 
applying the projection $\pi_{pl\rightarrow p}$.

For $k=11$, the lifting diagram is
\[S_1\rightarrow S_2\rightarrow S_4\rightarrow S_8\rightarrow S_{16} \xrightarrow{\pi} T_1\rightarrow T_3\rightarrow T_9.\]
Empirically, $T_3$ would usually be empty so there is no need to lift to $T_9$. 

For $k\in\{10,12\}$, the goal for the lifting phase is different. Instead of showing that the final set is empty,
it is enough to show that the only remaining speed set is $\onetok$, thus we might not need to lift with $c=k+1$.
We do not need to lift $\onetok$ as Lemma~\ref{lem:corlarge} asserts that the speed set is eventually $(k,p)$-proper.
for all $p\geq110$ (when $k=10$) and $p\geq156$ (when $k=12$), and Lemma~\ref{lem:subsetresolve} could be applied for smaller primes.

For both $k=10$ and $k=12$, the lifting phase is
\[S_1\rightarrow S_2\rightarrow S_4\rightarrow S_8 \xrightarrow{\pi} T_1\rightarrow T_3\rightarrow T_9\xrightarrow{\pi} U_1.\]

\section{Results}
\label{sec:result}

\begin{table}
\caption{Prime divisors of $u_1\cdots u_k$ for a counterexample $\ubf$. \label{table:primedivisors}}
\begin{tabular}{|c|l|c|c|}
\hline
$k$ & \multicolumn{1}{c|}{$P_k$}  
    & $\ln\displaystyle\prod_{p\in P_k}p$ & 
    \begin{tabular} [c]{@{}c@{}}$\ln B_k$\\ (Lemma~\ref{lem:proddivB}) \end{tabular} \\ \hline
10  & \begin{tabular}[c]{@{}l@{}}
      103, 131, 137, 139, 149, 151, 157, 107, 109, 127, 163, 167, 173, \\ 
      179, 433, 191, 193, 197, 199, 211, 223, 227, 229, 233, 239, 241, \\ 
      251, 257, 263, 269, 271, 277, 281, 283, 293, 307, 311, 313, 317, \\ 
      331, 337, 347, 349, 353, 359, 367, 373, 379, 383, 389, 401, 409,\\ 
      419, 421, 431, 433, 439, 443, 449, 457, 461
      \end{tabular}                                       
    & 339  & 337.7 \\ \hline
11  & \begin{tabular}[c]{@{}l@{}}
      23, 131, 137, 139, 149, 151, 157, 163, 167, 173, 179, 181, 191, \\ 
      193, 197, 199, 211, 223, 227, 229, 233, 239, 241, 251, 257, 263,\\ 
      269, 271, 277, 281, 283, 293, 307, 311, 313, 317, 331, 337, 347,\\ 
      349, 353, 359, 367, 373, 379, 383, 389, 397, 401, 409, 419, 421,\\ 
      431, 433, 439, 443, 449, 457, 461, 463, 467, 479, 487, 491, 499, \\
      503, 509, 521, 523, 541, 547, 557, 563, 569, 571, 577
      \end{tabular}                                                                              
    & 435 & 434.5 \\ \hline
12  & \begin{tabular}[c]{@{}l@{}}
    149, 151, 167, 179, 181, 191, 193, 197, 199, 211, 223, 227, 229, \\ 
    233, 239, 241, 251, 257, 263, 269, 271, 277, 281, 283, 293, 307, \\ 
    311, 313, 317, 331, 337, 347, 349, 353, 359, 367, 373, 379, 383, \\ 
    389, 397, 401, 409, 419, 421, 431, 433, 439, 443, 449, 457, 461, \\ 
    463, 467, 479, 487, 491, 499, 503, 509, 521, 523, 541, 547, 557, \\ 
    563, 569, 571, 577, 587, 593, 599, 601, 607, 613, 617, 619, 631, \\
    641, 643, 647, 653, 659, 661, 673, 677, 683, 691, 701, 709, 719, \\ 
    727 
      \end{tabular} & 550 & 545.3 \\ \hline
\end{tabular}
\end{table}

The prime divisors obtained from the algorithms in Section~\ref{sec:algo} are recorded in Table~\ref{table:primedivisors}. 
We note that listed primes are not necessarily minimal: they are the primes used in our verification, and a smaller set of primes, 
or even a less expensive set to compute, may exist.
We have also checked that every prime use for $k=10$ and $k=12$ satisfies the condition of Lemma~\ref{lem:subsetresolve}.
With that, we have verified that all $\vbf\in\Int{p,1}^k$ are eventually $(k,p)$-proper, 
for each $k \in \set{10, 11, 12}$ and $p\in P_k$. 

We are now ready to prove the main theorem.

\begin{theorem}
  The Lonely Runner Conjecture holds for $k\leq12$.
\end{theorem}
\begin{proof}
  Recall that $LRC(k)$ holds for $k\leq 9$ \cite{Trakulthongchai}. 
  From Corollary~\ref{cor:equivclasscheck} and the above computer verification, the product $u_1\cdots u_{10}$ of 
  any counterexample $\ubf$ to $LRC(10)$ must be divisible by $\prod_{p \in P_{10}}$. 
  This means $u_1\cdots u_{10} \geq \prod_{p \in P_{10}}p > B_{10}$, contradicting Lemma~\ref{lem:proddivB}.
  Hence no counterexample to $LRC(10)$ exists, i.e., $LRC(10)$ holds. 

  Similarly, as $u_1\cdots u_{k}\geq\prod_{p \in P_{k}}p>B_{k}$ for $k\in\{11,12\}$, $LRC(11)$ and subsequently $LRC(12)$ follow.
\end{proof}

\section{Implementation details}
\label{sec:implement}
Rosenfeld wrote \cite{rosenfeldcode} to verify the conjecture for $k=7$ and the second author developed \cite{tanupatcode}
to verify the case of $k=8,9$. The latter code takes 15 minutes for $k=8$ and approximately 23 hours for $k=9$. 
We developed algorithms in Section~\ref{sec:algo}, with some help of AI-based tools, 
to improve performance and scalability~\cite{mycode}. 

In addition to the mathematical improvements described in Sections~\ref{sec:quotient},~\ref{sec:sieves}~and~\ref{sec:onetok}, 
the practical aspect was also improved by refactoring the implementation to support parallel execution and
introducing more memory-efficient data structures.

One important note is that in the computation of $\tilde{I}(k, p, 1)$, the first coordinate is fixed to 1 and the remaining 
$k-1$ indices are filled with values from the range $[2, (p-1)/2]$. The ordering of elements in the remaining $k-1$ indices are arbitrary, 
chosen so that the search space can be pruned as much and as early as possible. In filling in the remaining $k-1$, we also skip over 
every tuple that corresponds to the same tuple under permutation, thus reducing the work by a factor of $(k-1)!$. By choosing that 
the first index is 1, the program would produce at most $k$ times as many sets in $\tilde{I}(k, p, 1)$.

We could not find a good way to reduce the search space down by another factor of $k$ despite the symmetry, 
so the congruence in the output sets collapse during the post-processing. In the accompanying log files, 
we show all speed sets that survived the lifts, before the post-processing, but note that this step can be done 
before the lifting step.

These refinements dramatically improve runtime. Compared to the second author's original implementation in \cite{Trakulthongchai}, 
we proved $k=8$ in 12 seconds and $k=9$ in 181 seconds on a 10-core Apple M4 processor. 
The proofs of $k=10$, $k=11$, and $k=12$ were run in multiple batches, with multiple versions, and on multiple machines,
 making the exact time taken incomparable to other cases. The source of the computation, along with the log files, 
can be found in the first author's GitHub repository \cite{mycode}.

We note here that the third phase of the computation for $k=10$ can be replaced by simply lifting normally with $c=11$. In fact, we 
almost completed the proof of $k=10$ without using Section~\ref{sec:onetok}. However, as lifting is especially expensive with 
large primes, the technique put forth by Section~\ref{sec:onetok} allows us to finish the proof of $k=12$.

% \section{Further work}
% \label{sec:further}
% TODO: wait until really done to see what's left

\bibliographystyle{amsplain}
\bibliography{ref}

\end{document}
